\chapter{Conclusiones Generales y Trabajo Futuro}
\label{cap:conclusiones}

% --- SECCIÓN: Síntesis del Trabajo ---
\section{Síntesis del Trabajo}
\label{sec:sintesis}
% ===================================================================

Esta tesis ha desarrollado un marco matemático completo para el
análisis de regímenes económicos basado en la \emph{Tasa de Cambio
Relativa Óptima Conmutada} (TCROC), desde sus fundamentos
algebraicos hasta su extensión no lineal mediante computación de
reservorio.  El trabajo se articuló en torno a un hilo conductor
progresivo: cada capítulo construyó sobre los resultados del
anterior, generando una jerarquía de modelos de complejidad
creciente con garantías teóricas formales y validación empírica
sobre datos reales del mercado hondureño de combustibles.

El bloque inicial de esta investigación estableció la base
del edificio teórico al formalizar la familia de operadores TCROC y
sus variantes (TCROCM, ETCROC y ETCROCM) como transformaciones
parametrizadas sobre series temporales.  Se demostraron las
relaciones de anidamiento entre variantes y se estableció que la
señal $\alpha_t$ admite interpretación como coeficiente de un modelo
AR(1) local estimado por mínimos cuadrados ponderados.  Esta
interpretación estadística, aparentemente elemental, resultó ser la
piedra angular de toda la construcción posterior: al dotar a
$\alpha_t$ de significado como \emph{tasa de cambio local}, se
habilitó su uso tanto como entrada para modelos estocásticos
discretos (cadenas de Markov) como para modelos dinámicos continuos
(reservorios recurrentes).

Las secciones anteriores extendieron el operador a una
formulación polinómica, demostrando teoremas de existencia y
unicidad del estimador bajo condiciones de no degeneración de los
datos en la ventana.  Estos resultados proporcionaron las garantías
analíticas necesarias para que la señal $\alpha_t$ estuviera bien
definida en cada instante temporal, condición indispensable para la
modelación Markoviana posterior.

Posteriormente, se formalizó la integración del
operador TCROC con cadenas de Markov de tiempo discreto, incluyendo
las propiedades de $\phi$-mixing, convergencia asintótica y la
conexión entre el gap espectral de la matriz de transición y la
velocidad de convergencia al equilibrio estacionario.  El resultado
central de este capítulo, según el cual la señal $\alpha_t$ discretizada
hereda propiedades de mezcla bajo condiciones verificables, sentó
las bases para la inferencia estadística en los capítulos aplicados.

El análisis sobre matrices de transición presentó la
metodología unificada de estimación de la matriz de transición
$\hat{\mat{P}}$ mediante mínimos cuadrados no negativos (NNLS),
incluyendo la prueba constructiva de existencia y unicidad de la
solución.  La formulación NNLS, que transforma el problema matricial
en un sistema lineal aumentado mediante el producto de Kronecker,
constituye una contribución metodológica central de la tesis: evita
la optimización no convexa de la verosimilitud que hace al MS-AR
clásico vulnerable a mínimos locales, garantizando convergencia global. Complementariamente, las probabilidades estocásticas de la matriz de transición ($\sum_i \hat{p}_{ij} = 1$ para cada columna $j$) se aseguran retrospectivamente mediante una normalización columna por columna, preservando la convención de matrices estocásticas por columnas adoptada en el Capítulo~\ref{cap:tcroc_hibrido_integrado}.

El Capítulo~\ref{cap:regimenes_combustibles} constituyó la
validación empírica integral del marco TCROC-Markov.  Aplicado a
series semanales de cuatro combustibles hondureños (Regular, Súper,
Diésel y Kerosene) durante 2017--2025, el modelo demostró
superioridad estadísticamente significativa sobre la discretización
por cuantiles fijos ($p < 0.05$ en exactitud para las cuatro series)
y estabilidad numérica total frente al benchmark MS-AR, que no
convergió en ningún escenario de validación.  El análisis espectral
reveló una estructura de mercado caracterizada por alta inercia
(gaps espectrales $\gamma \in [0.021, 0.040]$ y tiempos de mezcla de
25 a 47 semanas), formalizando cuantitativamente la ``viscosidad''
del mercado regulado hondureño.  Adicionalmente, se identificó y
formalizó el fenómeno de degeneración del modelo para ventanas $W$
excesivamente grandes, estableciendo la restricción $W \le 52$ como
condición necesaria para la identificabilidad de los regímenes.

El apartado teórico final cerró el arco analítico de la tesis
al demostrar que el marco TCROC-Markov es un caso particular de una
arquitectura más general: el Computador de Reservorio Estructurado
Estocástico (SSRC).  El Teorema de Inclusión (Teorema~6.2)
estableció formalmente que toda predicción del modelo Markoviano es
reproducible por un reservorio con $\mat{W}_{\text{res}} = \mat{0}$,
$\sigma = \text{id}$ y $a = 1$, pero no a la inversa.  La extensión
no lineal, validada empíricamente con mejoras del 20--23\% en RMSE
respecto al modelo Markoviano (estadísticamente significativas para
tres de cuatro series, $p < 0.05$ vía test de Diebold-Mariano),
demostró que la representación continua de estados mediante
reservorios recurrentes captura estructura que la discretización
Markoviana no puede representar, particularmente las asimetrías
temporales entre caídas y recuperaciones de precios.

% --- SECCIÓN: Resumen Cuantitativo ---
\section{Resumen Cuantitativo de Resultados}
\label{sec:resumen_cuantitativo}

La Tabla~\ref{tab:resumen_final} resume las métricas de rendimiento
del modelo TCROC-Markov para las cuatro series de combustibles
analizadas, junto con los métodos de referencia. Se reportan la
exactitud de predicción de régimen, el número de estados $K$, la
ventana óptima $W$ y la significancia estadística frente al
benchmark de cuantiles fijos.

\begin{table}[htbp]
\centering
\caption{Resumen consolidado de métricas por serie y método.}
\label{tab:resumen_final}
\small
\begin{tabular}{lccccc}
\hline
\textbf{Serie} & \textbf{$K$} & \textbf{$W$} &
\textbf{Exactitud} & \textbf{vs.~Cuantiles} &
\textbf{MS-AR} \\
\hline
Gasolina Regular & 3 & 12 & 91.2\% & $p < 0.01$ & No conv. \\
Gasolina Súper   & 3 & 12 & 89.7\% & $p < 0.01$ & No conv. \\
Diésel           & 4 & 8  & 87.4\% & $p < 0.05$ & No conv. \\
Kerosene         & 4 & 8  & 85.1\% & $p < 0.05$ & No conv. \\
\hline
\end{tabular}
\end{table}

En todos los casos, el modelo TCROC-Markov superó
significativamente al benchmark de cuantiles fijos.  El modelo
MS-AR no alcanzó convergencia en ningún escenario, lo que confirma
la ventaja de la formulación NNLS frente a la optimización por
máxima verosimilitud para estas series reguladas.

% --- SECCIÓN: Contribuciones ---
\section{Contribuciones}
\label{sec:contribuciones}
% ===================================================================

Las contribuciones de esta tesis se organizan en tres dimensiones:
teórica, metodológica y aplicada.

% --- SUBSECCIÓN: Contribuciones Teóricas --- - - - - - - - - - - - - - - - - - - - - - - - - - - - - - -
\subsection{Contribuciones Teóricas}

\begin{enumerate}
    \item \textbf{Jerarquía formal de operadores TCROC.}
    Formalización de las relaciones de anidamiento entre las
    variantes TCROC, TCROCM, ETCROC y ETCROCM, con demostración de
    que cada variante es un caso particular de la siguiente en la
    jerarquía.

    \item \textbf{Existencia y unicidad del estimador polinómico.}
    Demostración de que el operador TCROC polinómico admite un
    único estimador bajo condiciones de no degeneración, garantizando
    que la señal $\alpha_t$ está bien definida.

    \item \textbf{Propiedades de mezcla de la señal discretizada.}
    Prueba de que la secuencia de estados $\{S_t\}$ obtenida por
    discretización de $\alpha_t$ hereda propiedades de $\phi$-mixing,
    habilitando la inferencia estadística asintótica.

    \item \textbf{Existencia y unicidad de la solución NNLS.}
    Demostración constructiva de que el problema de estimación de
    $\hat{\mat{P}}$ mediante NNLS aumentado posee solución única
    cuando la matriz de diseño tiene rango completo, con
    caracterización de las condiciones de rango en términos de la
    representación de todos los estados en los datos.

    \item \textbf{Teorema de Inclusión TCROC-Markov $\subset$ SSRC.}
    Prueba formal de que el modelo Markoviano lineal es un caso
    degenerado del computador de reservorio estructurado,
    estableciendo una jerarquía estricta de expresividad entre ambos
    marcos.

    \item \textbf{Preservación de la Propiedad de Estado de Eco bajo
    integración de fuga.}  Extensión del teorema de ESP al caso
    Leaky-ESN, demostrando que la contracción del mapa de reservorio
    se preserva para todo $a \in (0, 1]$ cuando el radio espectral
    satisface $\rho(\mat{W}_{\text{res}}) < 1$.

    \item \textbf{Formalización del fenómeno de degeneración.}
    Caracterización analítica del colapso del modelo cuando $W \to T$,
    mostrando que $\text{Var}(\alpha_t) \to 0$ y K-Means converge a
    $K_{\text{eff}} = 1$ independientemente de $K$ nominal.
\end{enumerate}

% --- SUBSECCIÓN: Contribuciones Metodológicas --- - - - - - - - - - - - - - - - - - - - - - - - - - - - - - -
\subsection{Contribuciones Metodológicas}

\begin{enumerate}
    \item \textbf{Marco TCROC-Markov de dos etapas.}  Integración
    del operador TCROC con cadenas de Markov estimadas vía NNLS,
    que evita la optimización no convexa del MS-AR y garantiza
    convergencia global.  Este marco ha sido aceptado para
    publicación en IEEE CONCAPAN XLII (2025).

    \item \textbf{Protocolo jerárquico de optimización de
    hiperparámetros.}  Diseño de un criterio de selección en tres
    niveles (RMSE $\to$ Exactitud $\to$ AIC) con búsqueda exhaustiva
    en grilla y restricción fundamentada de $W \le 52$.

    \item \textbf{Arquitectura TCROC-SSRC.}  Extensión no lineal
    que reemplaza la discretización K-Means por un reservorio
    recurrente Leaky-ESN, preservando la estimación NNLS para la
    capa de lectura y habilitando representación continua de estados.

    \item \textbf{Pipeline de validación con test de
    Diebold-Mariano.}  Protocolo de evaluación estadística rigurosa
    que combina ventanas deslizantes, ensemble de 30 realizaciones y
    pruebas de significancia para la comparación de capacidad
    predictiva.
\end{enumerate}

% --- SUBSECCIÓN: Contribuciones Aplicadas --- - - - - - - - - - - - - - - - - - - - - - - - - - - - - - -
\subsection{Contribuciones Aplicadas}

\begin{enumerate}
    \item \textbf{Caracterización cuantitativa del mercado hondureño
    de combustibles.}  Identificación de 3 a 5 regímenes por
    combustible con centroides interpretables, distribuciones
    estacionarias y tiempos de mezcla que formalizan la inercia del
    mercado regulado.

    \item \textbf{Evidencia de superioridad sobre benchmarks
    establecidos.}  Demostración empírica de que TCROC-Markov supera
    a la discretización por cuantiles ($p < 0.05$) y que MS-AR no
    converge en este contexto, y de que TCROC-SSRC mejora sobre
    TCROC-Markov en 20--23\% ($p < 0.05$ para tres de cuatro
    series).

    \item \textbf{Descubrimiento de la jerarquía de complejidad de
    mercado.}  Hallazgo empírico de que Regular y Diésel ($K^* = 3$)
    tienen dinámica más simple que Súper ($K^* = 4$) y Kerosene
    ($K^* = 5$), con correlación positiva entre complejidad
    estructural y dificultad de predicción.
\end{enumerate}

% --- SECCIÓN: Limitaciones Generales ---
\section{Limitaciones Generales}
\label{sec:limitaciones_generales}
% ===================================================================

Es pertinente reconocer las limitaciones transversales del trabajo
para contextualizar el alcance de los resultados.

\textbf{Aplicación a una sola economía.}  Todos los resultados
empíricos se obtuvieron sobre datos del mercado hondureño de
combustibles.  Aunque el marco teórico es general, la validación
empírica está limitada a un mercado regulado con características
específicas (ajustes semanales, baja volatilidad relativa, cuatro
productos derivados del petróleo).  La generalización a mercados
desregulados o a otras clases de activos requiere validación
adicional.

\textbf{Horizonte de predicción unitario.}  Tanto el modelo
Markoviano como el SSRC operan con pronóstico a un paso adelante
($h = 1$ semana).  La extensión a horizontes múltiples, ya sea
mediante potencias de la matriz de transición o mediante esquemas
recursivos en el reservorio, no fue abordada y constituye trabajo
futuro.

\textbf{Condicionamiento numérico del reservorio.}  Como se
documentó en la sección de Reservorios Estocásticos, los números de
condición de la matriz de estados del reservorio alcanzan valores de
$\kappa(\mat{H}) \sim 10^{9}$--$10^{10}$, lo que indica sensibilidad
numérica a pesar de las garantías teóricas de unicidad.  Aunque la
restricción de no negatividad (NNLS) y el promediado de ensemble
mitigan este problema en la práctica, una solución formal,
por ejemplo, mediante el uso de la regularización de Tikhonov con selección
adaptativa del parámetro de regularización, queda pendiente.

\textbf{Supuesto de estacionariedad.}  El modelo Markoviano asume
probabilidades de transición constantes, y el SSRC asume pesos de
reservorio fijos tras el entrenamiento.  El periodo 2017--2025
incluye la pandemia de COVID-19 y el pico de precios de 2022, eventos
que potencialmente violan estos supuestos.  La evaluación formal de
la estabilidad temporal de los parámetros mediante partición de la
muestra o pruebas de cambio estructural no fue realizada.

% --- SECCIÓN: Líneas de Trabajo Futuro ---
\section{Líneas de Trabajo Futuro}
\label{sec:trabajo_futuro}
% ===================================================================

Los resultados y limitaciones de esta tesis sugieren varias
direcciones de investigación futura, organizadas por horizonte
temporal.

% --- SUBSECCIÓN: Extensiones Inmediatas --- - - - - - - - - - - - - - - - - - - - - - - - - - - - - - -
\subsection{Extensiones Inmediatas}

\textbf{Sistemas multivariados acoplados.}  La extensión natural del
marco es pasar de series univariantes a sistemas de series temporales
interrelacionadas.  En el contexto centroamericano, los precios de
combustibles de las seis economías del istmo están acoplados por
proximidad geográfica, acuerdos comerciales y dependencia compartida
del crudo WTI.  Un modelo TCROC-Markov multivariado estimaría una
matriz de transición conjunta $\hat{\mat{P}} \in
\mathbb{R}^{K^m \times K^m}$, donde $m$ es el número de series
acopladas, capturando las dependencias cruzadas entre mercados.
Esta línea de investigación se encuentra actualmente en desarrollo
como artículo de investigación independiente, con resultados
preliminares que sugieren la existencia de una estructura de tipo
\emph{Hub-and-Spoke} en la red centroamericana de precios.

\textbf{Pruebas de estabilidad temporal.}  La partición de la
muestra en subperiodos (pre-pandemia 2017--2019, pandemia 2020--2021,
post-pandemia 2022--2025) y la comparación de las matrices
$\hat{\mat{P}}$ estimadas en cada subperiodo, mediante pruebas de
homogeneidad basadas en distancia entre matrices estocásticas,
permitirían evaluar formalmente el supuesto de estacionariedad.

\textbf{Regularización adaptativa del SSRC.}  La sustitución de NNLS
por un esquema de regularización de Tikhonov con selección del
parámetro $\lambda$ por validación cruzada generalizada (GCV)
podría mejorar el condicionamiento numérico sin sacrificar la
interpretabilidad de los pesos no negativos, mediante una
formulación de NNLS regularizada
\[ \hat{\vect{w}} = \argmin_{\vect{w} \ge 0} \|\mat{H}\vect{w} - \vect{y}\|^2 + \lambda\|\vect{w}\|^2. \]

% --- SUBSECCIÓN: Extensiones de Mediano Plazo --- - - - - - - - - - - - - - - - - - - - - - - - - - - - - - -
\subsection{Extensiones de Mediano Plazo}

\textbf{Topología de reservorio estructurada.}  En el SSRC actual,
la matriz $\mat{W}_{\text{res}}$ es aleatoria y dispersa.  Trabajos
recientes en la literatura de reservoir computing sugieren que
topologías deterministas (e.g., anillo, ciclo con saltos) pueden
reducir la varianza entre realizaciones y mejorar la
reproducibilidad.  Explorar estas topologías en el contexto TCROC
es una extensión natural.

\textbf{Incorporación de variables exógenas.}  El marco actual es
puramente endógeno: $\alpha_t$ se calcula exclusivamente a partir
de los precios pasados.  La incorporación de variables exógenas,
tales como el precio WTI, el tipo de cambio o índices de actividad económica,
como entradas adicionales al reservorio (ampliando
$\mat{W}_{\text{in}}$ a múltiples canales) podría mejorar la
capacidad predictiva, especialmente en periodos de choques externos.

\textbf{Pronóstico multi-horizonte.}  La extensión del horizonte
de predicción de $h = 1$ a $h \in \{1, 4, 8, 12\}$ semanas, ya sea
mediante recursión del reservorio o mediante formulación directa
con múltiples salidas, permitiría evaluar la degradación de la
capacidad predictiva con el horizonte y determinar el alcance
práctico del marco para la planificación económica.

% --- SUBSECCIÓN: Visión de Largo Plazo --- - - - - - - - - - - - - - - - - - - - - - - - - - - - - - -
\subsection{Visión de Largo Plazo}

\textbf{Generalización a otras clases de activos.}  El marco
TCROC-Markov/SSRC no asume nada específico sobre la naturaleza de
la serie temporal subyacente.  Su aplicación a tasas de interés,
índices bursátiles, tipos de cambio o indicadores macroeconómicos
constituye un programa de investigación que podría validar la
generalidad del enfoque más allá del mercado de combustibles.

\textbf{Fundamentación teórica de la cota de perturbación.}  Las
cotas de perturbación derivadas previamente en el apartado de SSRC
son conservadoras por varios órdenes de magnitud.  El desarrollo de
cotas más ajustadas, posiblemente explotando la estructura específica
de la no negatividad de $\mat{W}_{\text{out}}$ y la contracción del
reservorio, es un problema abierto de interés tanto teórico como
práctico.

% --- SECCIÓN: Reflexión Final ---
\section{Reflexión Final}
\label{sec:reflexion_final}
% ===================================================================

El recorrido de esta tesis, que abarca desde la formalización de un operador
de tasa de cambio local hasta la demostración de que dicho operador,
integrado con cadenas de Markov, constituye un caso particular de un
computador de reservorio no lineal, ilustra una idea fundamental
en matemáticas aplicadas: la tensión productiva entre simplicidad
interpretable y complejidad expresiva.

El modelo TCROC-Markov, con sus $K$ estados discretos y su matriz
de transición de $K(K-1)$ parámetros libres, es transparente: un
economista puede leer $\hat{\mat{P}}$ y comprender la dinámica del
mercado.  El modelo TCROC-SSRC, con su reservorio de $D$ neuronas y
su representación continua de estados, es más poderoso, al ser entre un 20 y un 23\%
más preciso en el pronóstico, aunque resulte menos inmediatamente legible.
Ambos tienen su lugar.

Lo que unifica ambos extremos es la formulación NNLS, que garantiza
soluciones no negativas, convergencia global y, crucialmente, la
posibilidad de interpretar los pesos resultantes como
\emph{contribuciones} de cada componente al pronóstico.  Esta
propiedad, que podría parecer un detalle técnico, es en realidad lo
que hace que el marco completo sea útil más allá de la academia:
un modelo que mejora la predicción pero no puede explicar por qué
tiene valor limitado para la toma de decisiones en política
económica.

La aplicación a combustibles hondureños no fue elegida al azar.  En
una economía donde el precio del galón de diésel afecta directamente
el costo del transporte de alimentos, la capacidad de anticipar
regímenes de precios y de cuantificar la incertidumbre asociada
a dicha anticipación, tiene consecuencias prácticas tangibles.  Si
el marco desarrollado en estas páginas contribuye, aunque sea
marginalmente, a mejorar la planificación en ese contexto, habrá
cumplido su propósito.





% -------------------------------------------------------------------
% -------------------------------------------------------------------
% APÉNDICE A: Profundización Técnica del Modelo TCROC-Markov
% -------------------------------------------------------------------
% -------------------------------------------------------------------


